\chapter{كيس من القطع}
% full version: chapters/01b_neurontypes.tex

\pencilfig{1}{\pencilart{1}}{كيس من القطع: كلها تقريبًا زرقاء، وقليل منها أحمر، وحفنة من الأنواع الأخرى.}

تخيّل أنهم أعطوك كيسًا فيه ستة وثمانون مليار قطعة متشابهة المظهر، وطلبوا منك أن تفهم كيف رُكّب منها حاسوب. من أين ستبدأ؟ المهندس لن يتأمل شكل القطع. سيسأل: كم مدخلًا للقطعة، وكم مخرجًا، وماذا تُخرج.

العصبون، وهذا اسم تلك القطعة، بسيط البنية على نحو مدهش. مداخله بالآلاف: تصل إليه أسلاك من عصبونات أخرى. ومخرجه واحد. لكن هذا المخرج يتفرع كالشجرة، فتذهب النسخة نفسها من الإشارة إلى آلاف المستقبِلين. تخيّل شخصًا يستمع إلى ألف محدّث في آن واحد، ويجيب بجملة واحدة، لكن يسمعها ألف شخص آخر.

ما هذه الإشارة؟ نقرة قصيرة. يُطلق العصبون بين حين وآخر نبضة: قفزة في الجهد تدوم نحو جزء من ألف من الثانية. نقرات العصبون الواحد كلها متماثلة، كضربات طبل واحد. لذلك فكل ما يستطيع العصبون قوله مكتوب لا في علوّ النقرة، بل في \emph{توقيتها}: أكثر تواترًا أو أقل، فورًا أو بعد تأخير.

أما في داخله فيجري العصبون أبسط حساب يمكن تخيّله: يجمع الإشارات الواردة، فإذا تجاوز المجموع العتبة نقر هو نفسه. بعض المداخل يدفع المجموع إلى الأعلى، وبعضها إلى الأسفل.

\section*{كيف نفرز الكيس}

طرف سلك الخرج حين يبلغ العصبون المجاور لا يلمسه. تبقى بينهما فجوة ضيقة، وتنقل الكيمياء الإشارة عبرها: يقذف السلك رشة من مادة خاصة، هي الناقل العصبي، ويلتقطها الجار. ومن حسن الحظ أن كل عصبون تقريبًا يقذف المادة نفسها من جميع أطراف سلكه. إذن يمكن فرز العصبونات بحسب \emph{ما} تقذفه.

سنسمي نوع العصبون بالأحرف اللاتينية الأربعة الأولى من الاسم الإنجليزي لمادته. العصبون الذي يقذف الغلوتامات هو \nt{GLUT}. والذي يقذف حمض غاما-أمينوبيوتيريك هو \nt{GABA}. ثم تأتي \nt{DOPA} (الدوبامين)، و\nt{SERO} (السيروتونين)، و\nt{NORA} (النورأدرينالين)، و\nt{ACET} (الأسيتيل كولين)، وبضعة أنواع أخرى.

حين يُفرز الكيس تظهر صورة غير متكافئة إلى حد بعيد. من كل ألف عصبون، نحو تسعمئة وستين من نوع \nt{GLUT}. إشارتها تدفع المستقبِل نحو النقر: إنها ”زائد“. ونحو ستة وثلاثين من نوع \nt{GABA}، وإشارتها تبعده عن العتبة: إنها ”ناقص“. وفي القشرة نسبتها أعلى، من الخُمس إلى الربع. أما سائر الأنواع مجتمعة فبضعة عصبونات في كل مئة ألف. قطرة في بحر.

\section*{الهاتف والراديو}

لكن هذه القطرة مبنية على نحو مختلف تمامًا. سلك عصبون \nt{GLUT} أو \nt{GABA} يصل إلى جيران محددين، كخط هاتفي: لا يتلقى الإشارة إلا من يُتصل به. أما المجموعات الصغيرة من \nt{DOPA} و\nt{SERO} و\nt{NORA} و\nt{ACET} فتعمل كمحطات إذاعية. تتشعب أسلاكها في مناطق شاسعة من الدماغ، وكثيرًا ما تنسكب مادتها لا في الفجوة بل في الحيز المحيط، فيسمع الرسالة نفسها ملايين الخلايا دفعة واحدة.

في الخطوط الهاتفية تجري البيانات: هذا خط في الصورة، وذاك صوت بطبقة معينة. وفي الراديو تجري الإعدادات العامة: كم ينبغي أن يعلو صوت كل شيء، وهل يستحق الأمر التعلم، وهل حان وقت النوم. وهكذا في أي حاسوب: خلايا الذاكرة بالمليارات، وسجلات التحكم حفنة. ومع ذلك لا تعمل الآلة من دونها.

\section*{المعنى يختاره المستقبِل}

ثمة دقيقة هنا. المادة نفسها لا تعرف أهي ”زائد“ أم ”ناقص“. الجزيء مفتاح. وما يحدث حين يُدخَل المفتاح يقرره القفل: بروتين خاص على جهة المستقبِل. الدوبامين نفسه ينشّط بعض المستقبِلين ويخمد آخرين، بحسب القفل الذي لديهم. كأنه بث إذاعي واحد يفهمه المستمعون كلٌّ على طريقته.

\section*{إشارة ”أفضل مما كان متوقعًا“}

أشهر قناة إذاعية هي الدوبامين. ماذا ينقل؟ إليك تجربة. يُعطى حيوان قطرة عصير على غير توقع، فتستجيب عصبونات الدوبامين برشقة من النقرات. ثم يُضاء مصباح قبل العصير في كل مرة. يتعلم الحيوان أن المصباح يَعِد بالعصير، فتصبح الرشقة عند المصباح، ولا رشقة عند العصير نفسه. وإذا لم يُعطَ العصير بعد المصباح، \emph{تصمت} عصبونات الدوبامين في لحظة الترقب.

\pencilfig{101}{%
% three rows: a spike raster of a DOPA neuron; lamp (amber) and juice (blue drop) above the axis
\foreach \row/\y in {1/0,2/-1.05,3/-2.1}{
  \draw[pen] (0,\y) -- (6.2,\y);
  \foreach \s in {0.3,0.75,1.1,2.15,2.6,3.05,3.45,3.9,4.45,4.95,5.4,5.85}{
    \pgfmathtruncatemacro{\gap}{(\row==3)&&(\s>3.6)&&(\s<4.7)}
    \ifnum\gap=0 \draw[pcGray, line width=0.7pt] (\s,\y) -- (\s,\y+0.3);\fi}}
% row 1: juice without a lamp -> burst at the juice
\foreach \s in {4.0,4.08,4.16,4.24,4.33}{\draw[pcAmber, line width=0.8pt] (\s,0) -- (\s,0.45);}
\draw[dotpen=pcBlue, wash=pcBlue] (4.15,0.72) circle (0.09); \draw[pcBlue, line width=0.6pt] (4.07,0.76) -- (4.15,0.92) -- (4.23,0.76);
% row 2: lamp -> burst at the lamp, nothing at the promised juice
\foreach \y in {-1.05,-2.1}{
  \foreach \s in {1.5,1.58,1.66,1.74,1.83}{\draw[pcAmber, line width=0.8pt] (\s,\y) -- (\s,\y+0.45);}
  \draw[dotpen=pcAmber, wash=pcAmber] (1.65,\y+0.72) circle (0.1);
  \foreach \a in {30,90,150}{\draw[pcAmber, line width=0.5pt] ($(1.65,\y+0.72)+(\a:0.14)$) -- ($(1.65,\y+0.72)+(\a:0.24)$);}}
\draw[dotpen=pcBlue, wash=pcBlue] (4.15,-0.33) circle (0.09); \draw[pcBlue, line width=0.6pt] (4.07,-0.29) -- (4.15,-0.13) -- (4.23,-0.29);
% row 3: lamp, no juice -> a gap where the juice was expected
\draw[pcBlue, line width=0.6pt, densely dotted] (4.15,-1.38) circle (0.09); \draw[pcBlue, line width=0.6pt, densely dotted] (4.07,-1.34) -- (4.15,-1.18) -- (4.23,-1.34);
\draw[penlite=pcRed] (3.65,-2.22) -- (3.65,-2.28) -- (4.65,-2.28) -- (4.65,-2.22);
\node[lbl, text=pcRed, anchor=north] at (4.15,-2.3) {فجوة};
\node[lbl, anchor=east, align=right] at (-0.1,0.15) {عصير\\بلا مصباح};
\node[lbl, anchor=east, align=right] at (-0.1,-0.9) {مصباح،\\ثم عصير};
\node[lbl, anchor=east, align=right] at (-0.1,-1.95) {مصباح،\\ولا عصير};
}{رشقة عند عصير غير متوقع؛ ثم عند المصباح الذي يَعِد به؛ ثم فجوة في اللحظة التي لم يصل فيها العصير الموعود.}

القاعدة التي تفسر الحالات الثلاث: الدوبامين لا يقول ”جيد“، بل ”أفضل مما توقعت“. العصير غير المتوقع مفاجأة سارة. والموعود لا مفاجأة فيه. والموعود الذي لم يُعطَ مفاجأة مزعجة. وهذه الكمية بالضبط هي ما تستخدمه البرامج التي تتعلم من أخطائها. وسنعود إليها لاحقًا.

\section*{ما تحمله معك}

الدماغ شبكة هائلة من قطع متشابهة المظهر. كلها تقريبًا ”زوائد“ و”نواقص“ في شبكة البيانات الهاتفية. وقلة قليلة منها محطات إذاعية تضبط الشبكة كلها دفعة واحدة. في الفصل التالي سنرى ماذا يمكن أن نبني من ”الزوائد“ وحدها.

\fullversion{1}
